\documentclass[french, 11pt]{article}

\input{/home/theo/MP2I/setup.tex}

\def\chapitre{3}
\def\pagetitle{Calculabilité, décidabilité.}

\renewcommand*{\P}{\texttt{P}}
\newcommand*{\NP}{\texttt{NP}}

\begin{document}

\input{/home/theo/MP2I/title.tex}

\section{\texorpdfstring{$\P \et \NP$}{Lg}}

\begin{defi}{}{}
    $\P$ et $\NP$ sont des classes de complexité. Les classes de complexité sont des ensembles de problèmes que l'on rassemble en fonction de leur difficulté selon plusieurs critères.\\
    La classe $\P$ correspond aux problèmes de décision qui peuvent être résolus en temps polynomial.\\
    La classe $\NP$ correspond aux problèmes pour lesquels on peut vérifier une solution en temps polynomial.
\end{defi}

\vspace*{-0.6cm}

\begin{defi}{}{}
    Une instance $I$ d'un problème $P$ est dite positive si la réponse à la question est \texttt{Oui}.\\
    Un algorithme $V$ est un vérificateur de $P$ si $I$ positif $\iff \exists c$ tel que $V(I,c)=$ \texttt{Oui}. 
\end{defi}

\vspace*{-0.6cm}

\begin{ex}{}{}
    Étant donné $G$ un graphe et $k$ un entier, est-il possible de colorer $S$ avec les $k$ premiers entiers en évitant d'avoirs deux voisins de même couleur ?\\
    Un vérificateur : pour tout sommet on vérifie si ses voisins sont de couleur différente.
\end{ex}

\vspace*{-0.6cm}

\begin{defi}{}{}
    Un problème $A$ appartient à $\NP$ si :\\
    --- Il existe un vérificateur $V$ qui s'exécute en temps polynomial en la taille de ses entrées.\\
    --- Il existe un polynome $R$ tel que la taille d'un certificat est bornée par $R(|I|)$.
\end{defi}

\vspace*{-0.6cm}

\begin{ex}{$\NP$ ?}{}
    Problème $A$ : étant donnés un graphe $G$, un entier $k$, existe-t-il un sous-graphe $G'$ à $k$ sommets tels qu'ils soient tous voisons deux-à-deux ?\\
    Problème $B$ : étant donnés une formule $\phi$, est-elle satisfiable.
\end{ex}

\vspace*{-0.6cm}

\begin{thm}{$\P\subset\NP$}{}
    Soit $P_1\in \P$. Il existe $A$ un algorithme polynomial pour le résoudre.\\
    On construit $V_p$ prenant une instance $I$ de $P$, un certificat $C$ et renvoyant $A(I)$.\\
    Alors $V_p$ s'exécute bien en temps polynomial, et c'est bien un vérificateur.\\
    $I$ positive $\iff$ $A(I)=\texttt{Oui}\iff\exists c, ~ V_p(I,c)=\texttt{Oui}$.\\
    On a $R=1$ convient puisqu'on peut choisir $c=\e$, alors $\P_1\in\NP$ donc $\P\subset\NP$.
\end{thm}

\vspace*{-0.6cm}

\begin{defi}{Réduction.}{}
    Une réduction de $P_1$ vers $P_2$ est une fonction $f$ qui à une instance de $P_1$ associe une instance de $P_2$ et :\\
    --- $I$ est une instance positive de $P_1$ ssi $f(I)$ est une instance positive de $P_2$.\n
    On dit d'une réduction qu'elle est polynomiale si elle se calcule en temps polynomiale en la taille de son entrée.
\end{defi}

\vspace*{-0.6cm}

\begin{ex}{}{}
    Ensemble Indépendant :\\
    --- Instance : graphe $G$, entier $k$.\\
    Existe-t-il $S'\subset S$ de taille $k$ tel que $\forall s,t \in S', ~ (s,t)\notin A$.\n
    Couverture par les sommets :\\
    --- Instance : graphe $G$, entier $k$.\\
    Existe-t-il $S''\subset S$ tel que $\forall (x,y) \in A, ~ x \in S'' \ou y \in S''$ et $|S''|=k$.\n
    Ensemble Indépendant se réduit polynomialement au problème Couverture par les sommets.
    \tcblower
    \boxed{\ra} Soit $f:(G,k)\mapsto (G,|S|-k)$. Montrons que $f$ est une réduction.\\
    Soit $(G,k)$ positive de Ensemble indépendant. Il existe $S'$ qui convient.\\
    Alors $\forall s,t \in S', ~ (s,t) \notin A$, $\forall (s,t) \in A, ~ s\notin S' \ou t \notin S'$.\\
    Donc $\forall (s,t) \in A, ~ s \in S \setminus S' \ou t \in S \setminus S'$.\\
    $|S\setminus S'| = |S| - k$. Donc $|S''|=S\setminus S'$ convient.\\
    \boxed{\la} Soit $(G,k)$ instance de Ensemble indépendant telle que $f(G,k)$ est positive de Couverture par les sommets.\\
    Il existe $S''$ de taille $|S|-k$ tel que $\forall (x,y) \in A, ~ x\in S'' \ou y \in S''$.\\
    $\forall s,t \notin S'', ~ (s,t)\notin A$, $\forall s,t \in S\setminus S'', ~ (s,t)\notin A$. Donc $S\setminus S''$ est un ensemble indépendant de taille $k$ de $G$, $(G,k)$ est une instance positive de Ensemble Indépendant.
    La réduction est polynomiale. Ensemble indépendant se réduit polynomialement à Couverture par les sommets.
\end{ex}

\begin{prop}{}{}
    Soit $f$ une réduction polynomiale de $P_1$ vers $P_2$ et $g$ de $P_2$ vers $P_3$.\\
    Alors $g\circ f$ est une réduction polynomiale de $P_1$ vers $P_3$.
    \tcblower
    $I$ instance positive de $P_1$ $\iff$ $f(I)$ instance positive de $P_2$ $\iff$ $g\circ f(I)$ instance positive de $P_3$.\\
    On calcule $f(I)$ en temps polynomial en taille de $I$, $g\circ f(I)$ en temps polynomiale en taille de $f(I)$.
\end{prop}

\begin{prop}{}{}
    Si $f(P_1) \subset P_2$, et que $P_2$ est de classe $\P$, alors $P_1$ l'est aussi.
    \tcblower
    Il existe $A$ un algorithme résolvant $P_2$ en temps polynomial. Il existe $f$ une réduction polynomiale.\\
    On pose l'algorithme $B$ renvoyant $A(f(I))$.\\
    Alors $I$ positif $\iff$ $f(I)$ positif $\iff$ $A(f(I))$ vrai $\iff$ $B(I)$ vrai.\\
    Donc $B$ résout $P_1$ en temps polynomial. Donc $P_1$ est de classe $\P$.  
\end{prop}

\begin{prop}{}{}
    Si $P_1$ se réduit polynomialement à $P_2$, alors $P_2\in\NP\ra P_1\in\NP$.
    \tcblower
    Il existe un vérificateur $V$ et un polynôme $R$ qui conviennent et une réduction polynomiale $f$ de $P_1$ vers $P_2$.\\
    Vérificateur $V'$ : instance $I$, certificat $c$ $\mapsto$ $V(f(I),c)$.\\
    $I$ positive pour $P_1$ $\iff$ $f(I)$ est positive pour $P_2$ $\iff$ $\exists c, ~ |c|\leq R(|f(I)|)$ et $V(f(I),c)$ vrai\\
    $\iff$ $\exists c, ~ |c|\leq R(|f(I)|)$ et $V'(I,c)$ vrai $\iff \exists c, ~ |c|\leq R'(|I|)$ et $V'(I,c)$ vrai.\\
    Donc $P_1\in\NP$. 
\end{prop}

\begin{defi}{}{}
    Un problème $P$ est dit $\NP$-difficile si $\forall P'\in\NP, ~ P'$ se réduit polynomialement à $P$.\\
    Un problème $P$ est dit $\NP$-complet si il est $\NP$-difficile et dans $\NP$. 
\end{defi}

\vspace*{-0.5cm}

\begin{prop}{}{}
    Si $P_1$ est $\NP$-complet et se réduit polynomialement à $P_2$, alors $P_2$ est $\NP$-difficile.
    \tcblower
    Soit $P'\in\NP$. Il se réduit polynomialement à $P_1$, il se réduit polynomialement à $P_2$, donc $P_2$ est $\NP$-difficile.
\end{prop}

\vspace*{-0.5cm}

\begin{thm}{de Cook. (Admis)}{}
    Le problème SAT est $\NP$-difficile, il est donc $\NP$-complet.
\end{thm}

\vspace*{-0.5cm}

\begin{ex}{3-SAT}{}
    Le problème 3-SAT est $\NP$-complet.
    \tcblower
    Il est dans $\NP$, même vérificateur que pour 2-SAT.\\
    Soit $\phi$ une instance de SAT-FNC, conjonction de $n$ clauses $c_i$.\\
    Pour chaque $c_i$, on modifie en fonction de sa forme :\\
    --- $c_i = l_j \rightarrow l_j \lor l_j \lor l_j$.\\
    --- $c_i = l_j \lor l_k \rightarrow l_j \lor l_k \lor l_j$.\\
    --- $c_i = l_j \lor l_k \lor l_l \rightarrow l_j \lor l_k \lor l_k$.\\
    --- $c_i = l_{i_1} \lor ... \lor l_{i_n} \rightarrow (l_{i_1}\lor l_{i_2} \lor y_{i_1}) \land (\lnot y_{i_1} \lor l_{i_3} \lor y_{i_4}) ... \land (\lnot y_{i(n-3)} \lor l_{i(n-1)} \lor l_{i_n})$.\\
    On a défini $f$, on pose donc $\phi'$ la conjonction des $f(c_i)$.\\
    \boxed{\ra} Supposons $\phi$ satisfiable. Il existe $v$ valuation telle que $v(\phi)=\top$.\\
    Alors $\forall i \in \lb1,n\rb, ~ \exists l_i \in c_i, ~ v(l_i)=\top$, si $|c_i|\leq3$, alors $v(f(c_i))=\top$.\\
    Si $|c_i|\geq4$, il existe $l_{ij}$ tel que $v(l_{ij})=\top$.\\
    --- Si $j\leq2$, on fixe $v'(y_ik)=\bot$ pour tout $k\in\lb1,n-3\rb, ~ v'(f(c_i))=\top$.\\
    --- Si $j\geq n-1$, $\forall k \in \lb1,j-2\rb$, $v(y_{ik})=\top$, les autres à $\bot$.\\
    Pour chaque clause, on peut étendre $v$ en $v'$ satisfaisant l'image par $f$, donc $v'(\phi')=\top$.\\
    \boxed{\la} Soit $\phi'$ satisfiable, $\exists v', ~ v'(\phi')=\top$.\\
    On considère $v$ : $v'$ restreinte aux variables de $\phi$. Montrons $v(\phi)=\top$.\\
    --- $c_i$ est de taille $\leq 3$ : $v'(f(c_i))=\top \iff v(c_i) = \top$.\\
    --- $v'(f(c_i))=\top$ sinon.\\
    Il y a $n-2$ clauses dans $f(c_i)$ et $n-3$ nouvelles variables. Chaque variable satisfait au plus une clause, il y a une clause qui est satisfiable par un littéral :\\
    --- $v'(l_{ik})=\top \ra v(l_{ik})= \top \ra v(c_i) = \top$.\\
    Donc $\forall i, ~ v(c_i) = \top \ra v(\phi)=\top$, donc $\phi$ satisfiable.\\
    La fonction est bien une réduction de SAT-FNC à 3-SAT.\n
    La réduction est bien polynomiale : dans le cas d'une clause de taille supérieure à 4, on crée $n-2$ clauses de taille 3, chacune en temps polynomial.
\end{ex}

\begin{ex}{}{}
    Montrer que Clique est $\NP$-complet.
    \tcblower
    Soit une formule sous FNC. On fait le graphe des littéraux, avec des arêtes ssi ils ne sont pas dans la même clause.\\
    On cherche une clique dans le graphe obtenu, elle sera de taille au plus $K$, le nombre de clauses. Montrons que c'est une réduction.\\
    \boxed{\ra} Supposons $\phi$ satisfiable. Il existe $v$ tel que $v(\phi)=\top$. Pour chaque clause $c_i$, il y a au moins un $l_{ik}$ tel que $v(l_{ik})=\top$. On peut choisir un tel littéral pour chaque clause pour former un ensemble $S'$ de sommets.\\
    Alors $\forall l_{ik}, l_{jk} \in S', ~ l_{ik}\neq l_{jk} \ra i \neq j$ et $v(l_{ik})=v(l_{jk})=\top$.\\
    Ils ne sont pas négation l'un de l'autre et appartiennent à des clauses différentes, donc $(l_{ik}, l_{jk})\in A$ et $|S'|=K$.\\
    \boxed{\la} Supposons que $G$ admette une clique de taille $K$ : $\exists S' \subset S, ~ |S'|=K \et \forall s,t \in S', ~ s\neq t \ra (s,t)\in A$.
    On fixe $v$ une valuation telle que $\forall l_{ik}\in S', ~ v(l_{ik})=\top$, bien définie car $S'$ ne contient pas un élément et sa négation.
    Chaque clause a au plus un élément dans $S'$, puisque $|S'|=K$, chaque clause a exactement un élément dans $S'$.\\
    Donc chaque clause possède un élément évalué à $\top$ par $v$, donc chaque clause est évaluée à $\top$, donc $v(\phi)=\top$.\n
    C'est une réduction; elle est bien polynomiale car pour $n$ clauses, on a $3n$ sommets et au plus $9n(n-3)$ arêtes.
\end{ex}

\begin{ex}{}{}
    \bf{Cycle Hamiltonien} :\\
    Instance : Graphe $G$.\\
    Question : Existe-t-il un cycle qui passe une unique fois par chaque sommet ?\n
    \bf{Voyageur de commerce} :\\
    Instance : Graphe complet pondéré $G$, un entier $k$.\\
    Question : Existe-t-il un cycle qui passe une unique fois par chaque sommet de poids total inférieur à $k$ ?\\
    On admet Cycle Hamiltonien $\NP$-difficile, montrons que voyageur de commerce $\NP$-complet.
    \tcblower
    Voyageur de commerce est $\NP$. Vérificateur : instance I et liste de sommets :
    Tient une liste des sommets vus, et la somme des poids des arêtes vues. Renvoie Vrai si Vus ne contient pas faux et somme de poids inférieur à k.\n
    Réduction : il existe un cycle qui passe exactement une fois par tous les sommets, on fixe $b$ le poids des arêtes bleues à 1, et celui des rouges à $|S|+1$. Un tel cycle dans un tel graphe ne passe que par des arêtes bleues et une fois par chaque sommet a un poids $|S|$. Si le cycle emprunte une arête rouge, son poids est d'au moins $|S|+1$. On prend donc $K=|S|$. On a construit une réduction. En effet, si $G$ contenait un cycle hamiltonien, son image contient un cycle hamiltonien composé d'arêtes bleues, un cycle de poids total inférieur à $K$.\\
    Si l'image contient un cycle hamiltonien de poids total inférieur à $K$, il n'est composé que d'arêtes bleues et était donc présent dans $G$.
\end{ex}

\begin{ex}{Sac à dos.}{}
    Le problème appartient à la classe $\NP$, il est $\NP$-complet.\\
    On étudie un algorithme de programmation dynamique. On étudie $V_{i,p}$ la valeur maximale que l'on peut transporter en n'utilisant que les $i$ premiers objets et avec un poids total inférieur ou égal à $p$.\\
    On a alors $V_{i+1,p}=\max(V_{i,p}, V_{i,p-p_{i+1}} + v_{i+1})$, $V_{0,p}=0$.\\
    \begin{algorithm}[H]
        \caption{Algorithme de programmation dynamique sac à dos.}
        \Entree{Deux tableaux $V_i$ et $P_i$, un entier $P$ et un entier $V$.}
        $T$ $\gets$ Tableau de taille $(n+1)\times(p+1)$.\\
        \Pour{$p$ allant de 0 à $P$}{
            $T[0][p]=0$.
        }
        \Pour{$i$ allant de $1$ à $n$}{
            \Pour{$p$ allant de 0 à $p_i-1$}{
                $T[i][p]=T[i-1][p]$.
            }
            \Pour{$p$ allant de $p_i$ à $P$}{
                $T[i][p]=\max(T[i-1][p], T[i-1][p-P_i]+v_i)$.
            }
        }
        \Retour{$T[n][p]\geq v$}.
    \end{algorithm}
    Ce programme s'exécute en $O(nP)$ et décide bien si une instance de sac-à-dos est positive.\\
    On a l'impression d'avoir une solution polynomiale d'un algorithme $\NP$-complet, c'est une illusion.\\
    En effet, $P$ est de taille $O(\log P)$, donc la taille de l'instance : $O(\sum(\log v_i + \log p_i) + \log P + \log V)$.\\
    Donc $P=2^{\log P}$, donc $nP$ n'est pas polynomial en la taille de l'instance. Notre algorithme ne prouve pas que sac-à-dos appartient à $\P$.
\end{ex}

\begin{ex}{Bin Packing}{}
    Instance : $n$ nombres rationnels $p_i/q_i$, un entier $m$.\\
    Question : Existe-t-il une partition des $n$ nombres en $m$ ensembles $B_j$ tels que $\forall j, ~ \sum_{B_j}\frac{p_i}{q_i}\leq 1$ ?\\
    Vérificateur : Entrée : Instance, On crée $L$ une liste composée de $(j,\sum_{i\in B_j}\frac{p_i}{q_i})$.
    Pour $i$ 1 à $n$, si $(T[i],x)$ est dans $L$, le remplacer par $(T[i],x+\frac{p_i}{q_i})$.\\
    Sinon, ajouter $(T[i],\frac{p_i}{q_i})$. Renvoyer si tous les couples de $L$ ont un second élément $\leq 1$.\\
    Toutes les opérations sur celles-ci sont donc réalisées en temps polynomial en $n$.
\end{ex}

\begin{ex}{Partitionnement Carré}{}
    Instance : $n$ entiers $t_i$, un entier $N$.\\
    Question : Est-il possible de partitionner le carré de taille $N\times N$ en $n$ carrés de tailles $t_i\times t_i$ ?\\
    Vérificateur : Instance, Certificat : position des carrés à placer.\\
    On vérifie qu'il n'y a pas de superpositions, de débordements et que la somme des aires est égale à $N^2$.\n
    Variante : $n$ couples $(t_i, e_i)$, un entier $N$.\\
    Question : Est-il possible de partitionner le carré de taille $N\times N$ avec $e_i$ carrés de taille $t_i\times t_i$.\\
    La taille de l'instance est beaucoup plus compacte; ce qui rend l'ancien vérificateur non polynomial.
\end{ex}

\section{Décidabilité.}

\begin{defi}{}{}
    On dit qu'un problème de décision est décidable s'il existe un algorithme qui s'arrête sur toutes ses instances et renvoie Vrai ssi elle est positive.
\end{defi}

\begin{thm}{Machine universelle.}{}
    Il existe un algorithme, dit \bf{machine universelle}, qui étant donné une fonction $f$ et une entrée $x$, simule $f(x)$.
\end{thm}

\begin{ex}{Problème de l'Arrêt.}{}
    Instance : Une fonction $f$, un paramètre $x$.\\
    Question : La fonction $f$ se termine-t-elle sur l'entrée $x$ ?\\
    Ce problème est indécidable.
    \tcblower
    Supposons qu'un tel algorithme existe.\\
    On définit \texttt{Paradoxe} tel que :\\
    --- Si \texttt{Arrêt}($f, f$), alors boucle infinie.\\
    --- Sinon, renvoyer vrai.\\
    Alors \texttt{Paradoxe(Paradoxe)} s'arrête ssi \texttt{Arrêt(Paradoxe, Paradoxe)=Faux} par définition, c'est absurde.
\end{ex}

\begin{prop}{}{}
    S'il existe une réduction de $P_1$ à $P_2$ et que $P_2$ est décidable, alors $P_1$ l'est aussi.
\end{prop}

\end{document}
